\scanpage{71}
\mspage{- 25 -}

\[
\left.\begin{aligned}
H_i(0;\mu_1,\ldots,\mu_q)&=0\\
mH_i(\unsure{n|m};\mu_1,\ldots,\mu_q)&=\handwritten{\sum'_{n,(r_j),(\lambda_\ell)}}G_i(n;r_1,\ldots,r_p;\handwritten{\lambda}_1,\ldots,\handwritten{\lambda}_q)\mathbin{.}\\
&\quad\mathbin{.}\handwritten{\prod_{j=1}^{p}\prod_{k=1}^{r_j}}H_i(s_{kj};\handwritten{\nu}_{kj1},\ldots,\handwritten{\nu}_{kjq})
\end{aligned}\right\}
\quad\text{pour }m>0.\tag*{(17)}
\]

la sommation étant étendue aux indices satisfaisant aux mêmes conditions (16) que dans (15). La comparaison de (15) et (17) montre par récurrence sur le degré total que l'on a

\[
\left|h_i(m;\mu_1,\ldots,\mu_{\unsure{q|c}})\right|\leq H_i(m;\mu_1,\ldots,\mu_q)
\]

pour tout système d'indices ; en d'autres termes, les séries $(H_i(m;\mu_1,\ldots,\mu_q)t^mw_1^{\mu_1}\ldots w_q^{\mu_q})$ sont des \uline{majorantes} pour les séries $(h_i(m;\mu_1,\ldots,\mu_q)x^mz_1^{\mu_1}\ldots z_q^{\mu_q})$ respectivement $(1\leq i\leq p)$. Tout revient encore à prouver qu'on peut choisir les $G_i(n;r_1,\ldots,r_p;\handwritten{\lambda}_1,\ldots,\handwritten{\lambda}_q)$ de sorte que les séries majorantes $(H_i(m;\mu_1,\ldots,\mu_q)t^mw_1^{\mu_1}\ldots w_q^{\mu_q})$ soient sommables dans un voisinage de $0$.

Or, la série entière dont la somme est fonction des $p+q+1$ variables réelles $t$, $v_i$ $(1\leq i\leq p)$ et $w_{\handwritten{\ell}}$ $(1\leq\handwritten{\ell}\leq q)$

\[
\frac{M}{(1-\alpha t)(1-\alpha\handwritten{\sum_{i=1}^{p}}v_i)(1-\alpha\handwritten{\sum_{\ell=1}^{q}}w_{\handwritten{\ell}})}
\]

est une \deleted{am}\handwritten{ma}jorante de chacune des séries (14) dès que $M$ est assez grand et $\alpha$ assez petit. Tout revient à prouver qu'il existe $p$ séries entières sa ns terme indépendant de $t$, qui satisfont aux équations

\[
\frac{dv}{dt}\,k=\frac{M}{(1-\alpha t)(1-\alpha\handwritten{\sum_{i=1}^{p}}v_i)(1-\alpha\handwritten{\sum_{\ell=1}^{q}}w_{\handwritten{\ell}})}
\qquad(1\leq k\leq p)\tag*{(18)}
\]

On déduit aussitôt de (18) que si ce système admet une solution, tous les $v_k$ sont égaux à $v_1=v$ satisfaisant à l'équation différentielle
% archive sheet number: 70

